\chapter{Introducción}\label{cap:introduccion}

Citamos para temporalmente para evitar errores en la parte de bibliografía \cite{borrego2019}

\section{Contexto y objeto del proyecto}
Definición del problema 

Hablar un poco de lo que hablamos en la primera tutoría, el querer buscar una solución fácil para profesores y entretenida para los alumnos de evaluar el aprendizaje haciendo cuestionarios en clase 
También explicar la necesidad de añadirle una verificación para los alumnos

\section{Motivación}
Motivación personal y profesional para la realización del proyecto

Aquí tendría un poco más de dudas, me gustaría hablar un poco de la experiencia personal con los cuestionarios en clase, y de la motivación a que futuros alumnos se sirvan de una plataforma como esta, enfocándome en la importancia de la enseñanza y evaluación en el proceso de aprendizaje, y cómo una herramienta como Quizzie puede mejorar la experiencia educativa tanto para profesores como para estudiantes.

\section{Objetivos del proyecto}
Objetivos académicos y técnicos/profesionales

En objetivos académicos me gustaría decir que quería plasmar todo lo aprendido durante la carrera tanto en aspectos técnicos como en organización y buenas prácticas
En objetivos técnicos/profesionales me gustaría hablar de la motivación de aprender nuevas tecnologías, y de saber desenvolverme en un proyecto de desarrollo de software completo, desde la concepción hasta la implementación y evaluación

\section{Resumen del Stack tecnológico}
Pequeña descripción inicial (se ampliará en capítulos posteriores)

Tengo pensado decir que el backend se ha desarrollado con FastAPI y el frontend con React, un poco como pone en el readme del proyecto, pero sin entrar en más detalles, y en otros capítulos posteriores explicar por qué se han elegido, ventajas frente a otras opciones, y más herramientas que he usado como uv, hablar de la bbdd, etc

\section{Estructura de la memoria}
Guía rápida de la memoria

Pequeña introducción al contenido de la memoria
